% TEMPLATE-MILC-v3.tex - plantilla maestra UNICA de las guias MILC v3.
%
% La usan las cuatro familias de guias, cada una con su driver:
%   · Grados 6-11          -> scripts/build-guias-g11.py
%   · Bebras               -> scripts/build-guias-bebras.py
%   · Semillero            -> scripts/build-guias-semillero.py
%   · Territorio interior  -> scripts/build-guias-territorio-interior.py
%
% Antes habia tres copias casi identicas (~400 lineas iguales, 6 distintas):
% cada mejora habia que aplicarla tres veces, y en la practica no se hacia
% (el motor de recursos vivia solo en dos de ellas). Lo que varia entre
% familias esta parametrizado en 6 placeholders que cada driver llena
% (se nombran aqui SIN su sintaxis real: el builder reemplaza por texto plano,
% tambien dentro de los comentarios, y un valor multilinea rompe el preambulo):
%
%   HEADER_IZQ        encabezado izquierdo de pagina
%   HEADER_DER        encabezado derecho de pagina
%   PORTADA_ETIQUETA  rotulo sobre el numero grande (GRADO/MOMENTO/MODULO)
%   PORTADA_SERIE     linea de serie de la portada
%   PORTADA_ID        identificador de la guia en la portada
%   PORTADA_CIRCULO   el circulito de grado (°), o vacio si no aplica
%   PDF_TITULO        titulo en texto plano (sin \textbf ni comillas TeX) para
%                     los metadatos del PDF (/Title); lo produce lib_xelatex.texto_pdf
%
% Composicion (auditoria 2026-09-03): las bandas de estacion piden espacio
% (\Needspace*) y prohiben el salto justo despues (\nopagebreak), asi no quedan
% huerfanas al pie; las cajas largas (softbox, drawbox, infoband, puentebox) son
% `breakable`, asi una caja de media pagina ya no arrastra todo a la siguiente
% dejando la anterior al 40 %. Todo texto sobre mostaza o verde va en milcNegro
% (blanco sobre #E5B400 da 1,93:1 y sobre #58A923 2,95:1; ninguno pasa WCAG AA).
% El resto de claves las documenta content/guias/_SCHEMA.md. El script valida
% que no queden placeholders sin reemplazar antes de escribir el archivo final.

% Etiquetado PDF (latex-lab): /Lang, estructura y texto alternativo de las
% figuras, para lectores de pantalla. Debe ir ANTES de \documentclass.
\DocumentMetadata{lang=es-CO,tagging=on}
\documentclass[11pt,letterpaper]{article}
% headheight=24pt: la cabecera derecha lleva dos lineas (identidad de la guia
% y estacion actual). headsep se acorta en la misma medida para que el bloque
% de texto quede exactamente donde estaba con headheight=12pt.
\usepackage[margin=2.0cm,headheight=24pt,headsep=13pt]{geometry}
\usepackage[table]{xcolor}
\usepackage{fontspec}
\usepackage[spanish]{babel}
\usepackage{tikz}
% Librerías TikZ para los diagramas de las guías (bloque `recursos:`):
%  · babel  → babel en español vuelve activos caracteres como «>», y TikZ los usa
%             en sus opciones (>=stealth, ->). Sin esta librería, cualquier
%             diagrama con flechas revienta con «\language@active@arg> has an
%             extra }». Debe ir ANTES de que se dibuje nada.
%  · positioning        → sintaxis «below left=2pt and 3pt».
%  · decorations.pathreplacing → llaves ({brace}) para señalar rangos.
\usetikzlibrary{babel,positioning,decorations.pathreplacing}
\usepackage{graphicx}
% Circuitos: el trabajo de electrónica se hace hoy con micro:bit y sus sensores
% integrados (MakeCode, sin cableado), así que no se necesitan esquemáticos.
% Si algún día entran montajes con protoboard, basta descomentar esta línea:
% los diagramas .tex/.tikz declarados en `recursos:` ya se incrustan solos.
% \usepackage{circuitikz}
\usepackage[most]{tcolorbox}
\usepackage{tabularx}
\usepackage{array}
\usepackage{fancyhdr}
\usepackage{titlesec}
\usepackage[document]{ragged2e}  % bandera a la derecha en todo el cuerpo
\usepackage{enumitem}
\usepackage{microtype}
\usepackage{etoolbox}
\usepackage{needspace}
% hyperref al final del preambulo: metadatos del PDF (titulo, autor, idioma,
% familia), marcadores por estacion (\pdfbookmark) y enlaces sin recuadro.
\usepackage[hidelinks,
  pdftitle={Ansiedad: qué es y qué no, y dónde está la línea},
  pdfauthor={PhD. Álvaro Cárdenas Orozco},
  pdfsubject={Territorio interior · Educación socioemocional · Grado 10° · Momento 1 · MILC v3},
  pdflang={es-CO},
  bookmarksopen=true,bookmarksnumbered=false]{hyperref}

% Helvetica Neue no trae la flecha U+2192, asi que cada «→» escrito literalmente
% en el YAML se compilaba a NADA, en silencio (462 flechas en 61 guias: rutas del
% tipo «paleta → tipografia → lockup» salian rotas). Mapeamos el caracter a su
% version matematica, que si tiene glifo. Asi el autor puede seguir escribiendo
% «→» comodamente en el YAML.
\usepackage{newunicodechar}
\newunicodechar{→}{\ensuremath{\rightarrow}}
% Mismo problema con los simbolos de marca y vinetas: el «✓» de «una palomita ✓»
% salia como cuadrito vacio en el PDF de todas las guias que lo usaban.
\usepackage{pifont}
\newunicodechar{✓}{\ding{51}}
\newunicodechar{✗}{\ding{55}}
\newunicodechar{✔}{\ding{52}}
\newunicodechar{•}{\textbullet}
\newunicodechar{≥}{\ensuremath{\geq}}
\newunicodechar{≤}{\ensuremath{\leq}}

\setmainfont{Helvetica Neue}
\setsansfont{Helvetica Neue}
\setlength{\parindent}{0pt}
\setlength{\parskip}{6pt}
% Sin particion de palabras: en 6.º-7.º una palabra cortada por guion al final
% de linea («Comuni-cacion») frena la lectura mucho mas que una linea corta.
\hyphenpenalty=10000
\exhyphenpenalty=10000
\pretolerance=2000
\tolerance=3000
\setlength{\emergencystretch}{3em}
\linespread{1.10}
\renewcommand{\arraystretch}{1.25}
\setlist{leftmargin=5mm,itemsep=1mm,topsep=1mm}
\newcolumntype{Y}{>{\RaggedRight\arraybackslash}X}

\definecolor{milcMagenta}{HTML}{E6007E}
\definecolor{milcVino}{HTML}{5A0038}
\definecolor{milcTurquesa}{HTML}{0093A5}
\definecolor{milcVerde}{HTML}{58A923}
\definecolor{milcMostaza}{HTML}{E5B400}
\definecolor{milcOcre}{HTML}{B86600}
\definecolor{milcNegro}{HTML}{241816}
\definecolor{milcGris}{HTML}{F7F7F4}
\definecolor{milcLinea}{HTML}{E8E2DD}
\definecolor{milcGradePrimary}{HTML}{0D9488}
\definecolor{milcGradeDark}{HTML}{115E59}
\definecolor{milcGradeSoft}{HTML}{CCFBF1}
\definecolor{milcGradeAccent}{HTML}{CFE7FF}
\definecolor{milcGradeText}{HTML}{FFFFFF}
\color{milcNegro}

\pagestyle{fancy}
\fancyhf{}
% Cabecera de dos lineas: identidad de la guia arriba y, a la derecha y en
% gris, la estacion en curso (\markright desde \milcestacion). Antes de la
% primera estacion la segunda linea va vacia; el \strut mantiene la altura
% para que la linea de arriba no baile entre paginas.
\lhead{\small Territorio interior · Educación socioemocional\\[-1pt]{\footnotesize\strut}}
\rhead{\small Grado 10° · Momento 1 · MILC v3\\[-1pt]{\footnotesize\strut\color{milcNegro!65}\rightmark}}
\cfoot{\small\thepage}
\renewcommand{\headrulewidth}{0pt}

% Sobre mostaza (#E5B400) y verde (#58A923) el texto va en milcNegro; sobre el
% resto de la paleta, en blanco. Se usa dentro de fonttitle/fontupper, donde un
% \color posterior manda sobre coltitle.
\newcommand{\milctintasobre}[1]{%
  \ifboolexpr{test{\ifstrequal{#1}{milcMostaza}} or test{\ifstrequal{#1}{milcVerde}}}%
    {\color{milcNegro}}{\color{white}}}

\newtcolorbox{titlebox}[1]{
  enhanced,colback=#1,colframe=#1,arc=7mm,boxrule=0pt,
  left=7mm,right=7mm,top=3mm,bottom=3mm,
  before skip=8pt,after skip=8pt,
  fontupper=\sffamily\bfseries\Large\color{white}
}

\newtcolorbox{softbox}[3]{
  enhanced,breakable,pad at break=3mm,
  colback=#2,colframe=#1,borderline west={4pt}{0pt}{#1},
  arc=2mm,boxrule=.4pt,left=4mm,right=4mm,top=3mm,bottom=3mm,
  before skip=6pt,after skip=8pt,title={#3},coltitle=white,
  colbacktitle=#1,fonttitle=\sffamily\bfseries\milctintasobre{#1},fontupper=\RaggedRight,
  attach boxed title to top left={xshift=3mm,yshift=-2mm},
  boxed title style={arc=2mm, boxrule=0pt}
}

\newtcolorbox{drawbox}[2]{
  enhanced,breakable,pad at break=4mm,
  colback=white,colframe=#1,arc=4mm,boxrule=1pt,
  left=5mm,right=5mm,top=4mm,bottom=4mm,before skip=8pt,after skip=8pt,
  title={#2},coltitle=white,colbacktitle=#1,fonttitle=\sffamily\bfseries\milctintasobre{#1},
  fontupper=\RaggedRight,
  attach boxed title to top left={xshift=4mm,yshift=-2mm},
  boxed title style={arc=3mm, boxrule=0pt}
}

\newtcolorbox{infoband}[1]{
  enhanced,breakable,pad at break=2mm,colback=#1!8,colframe=#1!50,arc=2mm,boxrule=.5pt,
  left=4mm,right=4mm,top=2mm,bottom=2mm,before skip=4pt,after skip=4pt,
  fontupper=\small\RaggedRight
}

% Puente narrativo entre secciones (contrato editorial Conéctate).
% Da continuidad: "Acabas de X. Ahora vas a Y porque Z."
% Visualmente: banda izquierda en color del grado, texto en itálica pequeña.
\newtcolorbox{puentebox}{
  enhanced,breakable,pad at break=2mm,colback=milcGradeSoft,colframe=milcGradeDark,
  borderline west={3pt}{0pt}{milcGradeDark},
  arc=0mm,boxrule=0pt,left=5mm,right=4mm,top=2.5mm,bottom=2.5mm,
  before skip=4pt,after skip=8pt,
  fontupper=\itshape\small\color{milcNegro}\RaggedRight
}
% La flecha va en modo matematico a proposito: Helvetica Neue NO tiene el glifo
% U+2192, asi que \textrightarrow se compilaba a NADA («Missing character: There
% is no → in font Helvetica Neue») y el puente salia sin flecha. En modo
% matematico la flecha viene de la fuente matematica y si se dibuja.
\newcommand{\puente}[1]{%
  \begin{puentebox}%
    \textcolor{milcGradeDark}{$\rightarrow$}\hspace{2mm}#1%
  \end{puentebox}%
}

\newcommand{\linead}[1]{\noindent\color{milcLinea}\rule{#1}{0.5pt}\color{milcNegro}}
\newcommand{\respuesta}[1]{\par\vspace{2mm}\foreach \n in {1,...,#1}{\noindent\color{milcLinea}\rule{\linewidth}{0.5pt}\par\vspace{5mm}}\color{milcNegro}}
\newcommand{\checkbox}{\textcolor{milcGradeDark}{\fbox{\phantom{x}}}\hspace{5pt}}

% ─── Listas aireadas para las guías socioemocionales ────────────────────
% Componer las verificaciones como «\checkbox texto\\» deja el texto que salta
% de línea debajo del cuadrito y apelmaza el bloque. Estos entornos alinean la
% continuación y airean los ítems. Los usa Territorio Interior; cualquier
% familia puede hacerlo.
\newlist{checklista}{itemize}{1}
\setlist[checklista]{
  label=\textcolor{milcGradeDark}{\fbox{\phantom{x}}},
  leftmargin=9mm, labelsep=3.2mm, itemsep=2.4mm,
  topsep=2.5mm, parsep=0pt, partopsep=0pt, before=\RaggedRight
}
\newlist{pasoslista}{enumerate}{1}
\setlist[pasoslista]{
  label=\textcolor{milcGradeDark}{\sffamily\bfseries\arabic*.},
  leftmargin=8mm, labelsep=3mm, itemsep=2.8mm,
  topsep=1mm, parsep=0pt, partopsep=0pt, before=\RaggedRight
}

\newcommand{\phasecard}[4]{
\begin{tcolorbox}[enhanced,colback=#3,colframe=#2,arc=3mm,boxrule=.5pt,height=3.05cm,valign=center,halign=center,left=1.5mm,right=1.5mm,top=2mm,bottom=2mm]
{\sffamily\bfseries\fontsize{22}{24}\selectfont\textcolor{#2}{#1}}\par
{\sffamily\bfseries\small #4}
\end{tcolorbox}
}

% ── Piezas del rediseno 2026 (legibilidad 6.º-11.º) ───────────────────
% Banda de estacion: reemplaza el titlebox macizo. Titulo a la izquierda,
% dato practico (tiempo/modalidad) a la derecha, en una sola linea.
% Nunca huerfana: pide 9 lineas libres (o salta de pagina rellenando con
% \vfill) y prohibe el salto justo despues de la banda. Nueve y no seis:
% la caja partible que sigue necesita ~80pt (titulo adosado, relleno y dos
% lineas) para arrancar en la misma pagina; con menos, tcolorbox la manda
% entera a la siguiente y la banda se queda sola. El penalty 10000 va
% en `after`, ANTES del espacio posterior: un `after skip` (aunque sea 0pt)
% es un \vskip tras la caja, y ese glue es un punto de corte legal que un
% \nopagebreak escrito despues de \end{tcolorbox} ya no cierra. Cada banda
% deja marcador PDF y actualiza la cabecera derecha.
\newcounter{milcestacion}
\newcommand{\milcestacion}[2]{%
\Needspace*{9\baselineskip}%
\stepcounter{milcestacion}%
\pdfbookmark[1]{#2}{milcestacion\themilcestacion}%
\markright{#2}%
\begin{tcolorbox}[enhanced,colback=#1,colframe=#1,arc=2mm,boxrule=0pt,
  left=5mm,right=5mm,top=2.6mm,bottom=2.6mm,before skip=11pt,
  after={\par\nopagebreak\vskip7pt}]
{\sffamily\bfseries\fontsize{15}{18}\selectfont\milctintasobre{#1}#2}
\end{tcolorbox}}

% Tarjeta de paso numerada: sustituye las tablas «Pilar / Aplicacion», que
% eran formato de consulta docente, no de estudiante. El numero va en un
% circulo sobre el borde para que la secuencia se lea de un vistazo.
\newcommand{\milcpaso}[3]{%
\Needspace{4\baselineskip}%
\begin{tcolorbox}[enhanced,colback=white,colframe=milcLinea,arc=1.5mm,boxrule=.9pt,
  left=14mm,right=4mm,top=3mm,bottom=3mm,before skip=4pt,after skip=4pt,
  fontupper=\RaggedRight,
  overlay={\node[anchor=north west,circle,fill=#2,text=white,inner sep=0pt,
    minimum size=7.4mm,font=\sffamily\bfseries\small]
    at ([xshift=3.6mm,yshift=-3.2mm]frame.north west) {#1};}]
#3
\end{tcolorbox}}

% Casilla marcable de tamano util con lapiz (la anterior era \fbox{\phantom{x}},
% de alto variable segun el texto que la seguia).
\newcommand{\milcmarca}{\framebox[3.8mm]{\rule{0pt}{3.4mm}}}

% Fila marcable de lista suelta (checks de la estacion 1).
\newcommand{\milccheckrow}[1]{%
\par\vspace{1.8mm}\noindent
\makebox[6.4mm][l]{\raisebox{-.55mm}{\textcolor{milcNegro}{\milcmarca}}}%
\parbox[t]{\dimexpr\linewidth-6.4mm\relax}{\RaggedRight #1}\par}

% Celda marcable dentro de la rubrica (tabla de autoevaluacion). Misma
% sangria francesa, pero pensada para vivir dentro de una columna X.
\newcommand{\milcopcion}[1]{%
\makebox[5.2mm][l]{\raisebox{-.55mm}{\textcolor{milcNegro}{\milcmarca}}}%
\parbox[t]{\dimexpr\linewidth-5.2mm\relax}{\RaggedRight #1}}

% ── Recursos visuales de la guía ──────────────────────────────────────
% Declarados en el bloque `recursos:` del YAML y emitidos por el builder.
% \guiaFigura   → imagen rasterizada o PDF (foto, carta celeste, montaje).
% \guiaDiagrama → fuente TikZ/circuitikz (.tex) incrustada como vector:
%                 es la vía para circuitos y esquemáticos editables.
% [#1] = texto alternativo (alt) · #2 = ruta relativa al .tex · #3 = pie
% El alt llega al PDF etiquetado (/Alt) y lo lee el lector de pantalla.
\newcommand{\guiaFigura}[3][]{%
\begin{center}
\includegraphics[width=0.88\linewidth,height=0.38\textheight,keepaspectratio,alt={#1}]{#2}\\[1.5mm]
{\footnotesize\itshape\textcolor{milcNegro!75}{#3}}
\end{center}
}

\newcommand{\guiaDiagrama}[2]{%
\begin{center}
\input{#1}\\[1.5mm]
{\footnotesize\itshape\textcolor{milcNegro!75}{#2}}
\end{center}
}

\begin{document}
\thispagestyle{empty}

% ── Portada ───────────────────────────────────────────────────────────
% Antes: un rectangulo de color a pagina completa. Se veia bien en pantalla,
% pero en la fotocopiadora del colegio se comia el toner y salia gris sucio.
% Ahora la banda ocupa el tercio superior y el resto va en blanco.
\begin{tikzpicture}[remember picture,overlay]
  \path[fill=milcGradePrimary]
    (current page.north west) rectangle ([yshift=-10.6cm]current page.north east);

  \node[anchor=north west,text=milcGradeText] at ([xshift=2.0cm,yshift=-1.55cm]current page.north west)
    {\sffamily\bfseries\fontsize{12}{14}\selectfont MOMENTO};
  \node[anchor=north west,text=milcGradeText,opacity=.85] at ([xshift=2.0cm,yshift=-2.35cm]current page.north west)
    {\sffamily\bfseries\fontsize{8.5}{10}\selectfont TERRITORIO INTERIOR · SOCIOEMOCIONAL};
  \node[anchor=north east,text=milcGradeText,opacity=.85,align=right] at ([xshift=-2.0cm,yshift=-1.55cm]current page.north east)
    {\sffamily\bfseries\fontsize{9}{12}\selectfont I.E. Sor María Juliana\\Cartago, Valle del Cauca};

  \node[anchor=west,text=milcGradeText] (gradonum) at ([xshift=1.85cm,yshift=-7.1cm]current page.north west)
    {\sffamily\bfseries\fontsize{112}{112}\selectfont 1};
  % El circulito de grado (°) solo aplica a las guias de grado («11°»); en
  % Bebras y Semillero el numero grande es un momento/modulo, no un grado.
  % Se posiciona relativo al nodo `gradonum`, no en coordenadas absolutas.
  

  \node[anchor=west,text=milcGradeText,text width=11.4cm,align=left] at ([xshift=1.5cm,yshift=0.5cm]gradonum.east)
    {
     {\sffamily\bfseries\fontsize{9}{11}\selectfont\MakeUppercase{Territorio interior · Socioemocional}}\\[3mm]
     {\sffamily\bfseries\fontsize{21}{25}\selectfont\setlength{\baselineskip}{27pt}Ansiedad\\[4pt]qué es y qué no\\[4pt]la que avisa y la que incapacita}\\[3.5mm]
     {\sffamily\bfseries\fontsize{15}{18}\selectfont Grado 10° · Momento 1}
    };
\end{tikzpicture}

\vspace*{9.4cm}
\vfill

{\sffamily\bfseries\fontsize{8}{10}\selectfont\textcolor{milcGradeDark}{LO QUE TE LLEVAS HOY}}

\begin{tcolorbox}[enhanced,colback=white,colframe=milcNegro,arc=2mm,boxrule=1.6pt,
  left=5mm,right=5mm,top=4mm,bottom=4mm,before skip=3pt,after skip=12pt,
  fontupper=\RaggedRight\fontsize{11}{15.5}\selectfont]
Tu mapa de la ansiedad: las señales que reconoces en tu propio cuerpo, tus cuatro situaciones pasadas por los criterios de interferencia, persistencia e intensidad, la lista de lo que has dejado de hacer por evitarla, y un acercamiento pequeño planeado con día.
\end{tcolorbox}

\begin{tcolorbox}[enhanced,colback=milcGris,colframe=milcGris,arc=2mm,boxrule=0pt,
  left=5mm,right=5mm,top=3.5mm,bottom=3.5mm,after skip=12pt]
{\sffamily\bfseries\fontsize{8}{10}\selectfont\textcolor{milcNegro!55}{ESTA GUÍA ES DE}}\\[3mm]
\begin{tabularx}{\linewidth}{X p{3.2cm} p{3.2cm}}
\linead{\linewidth} & \linead{3.0cm} & \linead{3.0cm}\\[-1mm]
{\fontsize{8.5}{10}\selectfont\textcolor{milcNegro!55}{Nombre completo}} &
{\fontsize{8.5}{10}\selectfont\textcolor{milcNegro!55}{Grupo}} &
{\fontsize{8.5}{10}\selectfont\textcolor{milcNegro!55}{Fecha}}\\
\end{tabularx}
\end{tcolorbox}

\vfill

\noindent\textcolor{milcLinea}{\rule{\linewidth}{1pt}}\\[2mm]
{\sffamily\bfseries\fontsize{9.5}{12}\selectfont Autor: Dr. Álvaro Cárdenas Orozco}\\[1mm]
{\sffamily\fontsize{8.5}{11}\selectfont\textcolor{milcNegro!55}{Metodología MILC · Apertura ancestral · Ansiedad · Triángulo Dussel-Estoicismo-Floridi}}

\newpage

\section*{Apertura: Ansiedad: qué es y qué no, y dónde está la línea}

\begin{softbox}{milcGradeDark}{milcGradeSoft}{Saber ancestral}
En el campo del Valle del Cauca hay un nombre para el sufrimiento que se siente \textbf{en el cuerpo}: la \textbf{«enfermedad de los nervios»}. Investigadoras la documentaron en \textbf{Sevilla y Barragán (Tuluá)}, y encontraron que la gente la describe con precisión física ---el pecho apretado, el temblor, la cabeza que no descansa, el cuerpo que no obedece--- \textbf{sin necesidad de un diagnóstico médico para nombrarla} (Piedrahita Forero y Tabares, 2021). \textbf{Y hay algo que esta guía no va a folclorizar, porque sería faltar a la verdad:} \emph{ese síndrome está documentado como \textbf{consecuencia de las violencias del conflicto armado}}. \textbf{No es una curiosidad del habla campesina: es la manera en que comunidades enteras nombraron lo que les dejó la guerra}, cuando no había psicólogo, ni palabra técnica, ni nadie preguntando. \textbf{Guarda las dos cosas que enseña.} \textbf{Primera:} \emph{la gente supo siempre que el sufrimiento se siente en el cuerpo} ---eso que la medicina tardó en aceptar, aquí ya tenía nombre---. \textbf{Segunda:} \emph{nombrarlo servía}. \textbf{Decir «estoy de los nervios» permitía pedir ayuda, que la familia entendiera y que alguien lo acompañara.} \emph{Un sufrimiento sin nombre no se puede contar; y lo que no se puede contar se aguanta solo.}
\end{softbox}

\begin{softbox}{milcTurquesa}{milcGris}{Saber contemporáneo}
¿Entonces sentir ansiedad es estar enfermo? \textbf{No, y confundir eso hace daño en las dos direcciones.} \emph{La ansiedad, en su forma corriente, \textbf{sirve}:} moviliza a prepararse, a practicar, a ensayar, y enciende una precaución razonable ante lo que sí es peligroso. \textbf{Sin nada de ansiedad nadie estudiaría para un examen ni miraría antes de cruzar.} \textbf{La línea que la separa de un trastorno no está en \emph{sentirla} sino en cuatro cosas concretas:} \textbf{cuánto interfiere} con la vida diaria, \textbf{cuánto dura} en el tiempo, \textbf{qué tan intensa} es, y \textbf{cuánto malestar} produce. \emph{Cuando causa disfunción y sufrimiento excesivos, deja de ser una señal útil y pasa a ser un problema que necesita atención.} \textbf{Y hay un mecanismo que conviene entender antes que cualquier otra cosa, porque explica por qué la ansiedad crece:} \emph{la \textbf{evitación} da alivio inmediato \textbf{y a mediano plazo la mantiene y la agrava}}, porque impide comprobar que lo temido no ocurre ---o que se puede tolerar---. \textbf{Por eso los tratamientos con evidencia hacen lo contrario de lo que pide el cuerpo:} \emph{un acercamiento \textbf{planificado y progresivo} a lo que se está evitando.}
\end{softbox}

\begin{softbox}{milcMostaza}{milcGris}{Pregunta-puente}
\textit{En Sevilla y en Barragán, ponerle nombre a lo que se sentía \textbf{permitía pedir ayuda}. \textbf{¿Qué has dejado de hacer este año porque te daba ansiedad} \ldots \textbf{y a quién no se lo has contado}?}
\end{softbox}

\begin{softbox}{milcVerde}{milcGris}{Saber y saber hacer}
Al terminar podrás: (1) \textbf{distinguir} la ansiedad que avisa de la que incapacita, usando interferencia, persistencia, intensidad y malestar; (2) \textbf{reconocer} tus propias señales corporales y nombrarlas; (3) \textbf{entender} por qué evitar alivia hoy y agrava mañana; (4) \textbf{planear} un acercamiento pequeño y concreto a algo que estás evitando ---y saber cuándo esto ya no se maneja solo---.
\end{softbox}



\small
\begin{tabularx}{\textwidth}{>{\bfseries}p{3.0cm}Y}
\rowcolor{milcGradePrimary!12}Autor & Dr. Álvaro Cárdenas Orozco\\
DBA & Comprende los fenómenos que afectan a su territorio, regula sus emociones ante ellos y acompaña a otros con empatía (transversal 6°--11°, articulado con la política «Escuela, territorio de vida» del MEN, Resolución 006519 de 2025, y la Ley 1620 de 2013)\\
Referentes & Primera ayuda psicológica (OMS, 2011/2012) · Directrices IASC (2007) · USGS y Servicio Geológico Colombiano · Pueblo nasa y la minga (CRIC) · Dussel · Estoicismo · Floridi\\
Producto final & Tu mapa de la ansiedad: las señales que reconoces en tu propio cuerpo, tus cuatro situaciones pasadas por los criterios de interferencia, persistencia e intensidad, la lista de lo que has dejado de hacer por evitarla, y un acercamiento pequeño planeado con día.\\
\end{tabularx}
\normalsize

\puente{Este es el primer momento del bloque más difícil del programa. \textbf{Abre nombrando lo que se siente en el cuerpo,} porque sin nombre no se puede pedir ayuda.}

\milcestacion{milcGradeDark}{Ruta MILC: así vamos a aprender}
\begin{tabularx}{\textwidth}{>{\centering\arraybackslash}X>{\centering\arraybackslash}X>{\centering\arraybackslash}X>{\centering\arraybackslash}X}
\phasecard{1}{milcVerde}{milcGris}{Escucha\\[-1mm]\small Observo, escucho y pregunto.} &
\phasecard{2}{milcTurquesa}{milcGris}{Sistematización\\[-1mm]\small La nombro y no la evito.} &
\phasecard{3}{milcMagenta}{milcGris}{Praxis\\[-1mm]\small Construyo y aplico.} &
\phasecard{4}{milcVino}{milcGris}{Evaluación\\[-1mm]\small Reflexión liberadora.}
\end{tabularx}


\puente{Primero, tu cuerpo: dónde se te siente. \emph{Sin diagnosticarte nada ---solo describir---.}}

\milcestacion{milcVerde}{Estación 1 de 3 · Escucha}

\begin{softbox}{milcVerde}{milcGris}{Escena para escuchar}
\textbf{Actividad 1 · IDENTIFICA --- Dónde se me siente.} \textbf{Nadie va a leer esta página.} \textbf{Primera parte.} Dibuja una silueta y \textbf{marca dónde se te siente} cuando estás nervioso: el pecho, el estómago, la garganta, las manos, la mandíbula, la cabeza, las piernas. \emph{Ponle palabras a cada marca ---«apretado», «revuelto», «temblor», «no me sale la voz»---.} \textbf{Segunda parte.} Escribe \textbf{cuatro situaciones} de este año en que sentiste eso. \textbf{Al frente de cada una, tres datos:} \emph{¿cuánto duró? ¿qué tan fuerte fue, de 1 a 10? ¿te impidió hacer algo?} \textbf{Tercera parte, la que importa.} Haz una lista de \textbf{lo que has dejado de hacer} para no sentir eso: hablar en público, ir a un lugar, preguntar en clase, llamar por teléfono, salir, presentarte a algo. \textbf{Y la última:} ¿alguien en tu casa \textbf{sabe} que te pasa esto? \emph{Si la respuesta es no, escribe por qué.}
\end{softbox}

{\sffamily\bfseries\fontsize{8}{10}\selectfont\textcolor{milcNegro!55}{MARCA CUANDO LO HAYAS HECHO}}
\milccheckrow{Marcaste dónde se te siente en el cuerpo, con palabras propias.}
\milccheckrow{Anotaste duración, intensidad y si te impidió hacer algo en las cuatro situaciones.}
\milccheckrow{Escribiste la lista de \textbf{lo que has dejado de hacer}.}

\begin{infoband}{milcVerde}


\textbf{En tu cuaderno (Actividad 1 · IDENTIFICA).} Título: «Actividad 1. Dónde se me siente». \textbf{Formato:} la silueta marcada + las cuatro situaciones con sus tres datos + la lista de lo evitado + si alguien lo sabe. \textbf{Extensión:} media página. \textbf{Sabes que terminaste cuando} tienes escrita \textbf{la lista de lo que evitas}. Esta página es tuya: \textbf{nadie más tiene que leerla si tú no quieres}.
\end{infoband}

\puente{Ya lo tienes. Ahora, para qué sirve la ansiedad, dónde está la línea y por qué evitar la hace crecer.}

\milcestacion{milcTurquesa}{Fase 2 · Sistematización: entender la ansiedad}

\begin{softbox}{milcTurquesa}{milcGris}{Cuatro claves sobre la ansiedad}
Cuatro claves para no confundir una señal útil con una enfermedad, ni una enfermedad con una debilidad.
\end{softbox}

\milcpaso{1}{milcTurquesa}{\textbf{La ansiedad sirve para algo, y por eso existe.} \emph{Moviliza a prepararse, a ensayar, a practicar; enciende la precaución ante lo que sí es peligroso.} \textbf{Sin ella nadie estudiaría para un examen ni miraría antes de cruzar una calle.} \emph{Y esto hay que decirlo de entrada, porque a los quince años circula la idea contraria:} \textbf{sentir ansiedad no significa estar enfermo}, igual que sentir tristeza no es tener depresión. \emph{El problema aparece cuando alguien se asusta de su propia ansiedad ---«¿me estará dando algo?»---,} \textbf{porque asustarse de una sensación la aumenta, y ahí empieza un círculo}.}
\milcpaso{2}{milcTurquesa}{\textbf{Dónde está la línea: cuatro criterios, no una intuición.} No se trata de qué tan feo se sienta, sino de: \textbf{(1) interferencia} ---¿te impide estudiar, dormir, salir, estar con gente?---; \textbf{(2) persistencia} ---¿lleva semanas o meses, no días?---; \textbf{(3) intensidad} ---¿es desproporcionada frente a lo que la dispara?---; y \textbf{(4) malestar} ---¿estás sufriendo de verdad?---. \emph{Cuando la ansiedad produce disfunción y sufrimiento excesivos, ya no es una señal: es un problema que necesita atención.} \textbf{Y fíjate en lo que esto te da:} \emph{un criterio para decidir si pedir ayuda, sin tener que diagnosticarte a ti mismo ---que no es tu trabajo ni el de esta guía---.}}
\milcpaso{3}{milcTurquesa}{\textbf{Evitar alivia hoy y la alimenta mañana.} Este es el mecanismo central y conviene entenderlo bien. \emph{Cuando algo te angustia y lo evitas, el alivio es inmediato ---y por eso el cerebro aprende que evitar funciona---.} \textbf{Pero a mediano plazo la evitación \emph{mantiene y agrava} el problema}, porque \emph{impide comprobar que lo temido no ocurre, o que sí se puede tolerar}. \textbf{Y hay un efecto adicional:} \emph{si cada señal de ansiedad lleva a escapar, las sensaciones del cuerpo empiezan a leerse como insoportables}, cuando en realidad fluctúan y bajan solas. \textbf{Por eso los tratamientos con evidencia hacen lo contrario de lo que pide el cuerpo:} \emph{acercarse, de forma planificada y progresiva, a lo que se está evitando.} \textbf{Nunca de golpe: por escalones.}}
\milcpaso{4}{milcTurquesa}{\textbf{El cuerpo también la nombra, y eso no es ignorancia.} \emph{Durante mucho tiempo se despreció que la gente dijera «es de los nervios» en vez de usar una palabra técnica.} \textbf{Y resulta que la descripción corporal era exacta:} el sufrimiento psíquico \textbf{se siente en el cuerpo}, y quien lo nombraba así estaba describiendo bien lo que le pasaba (Piedrahita Forero y Tabares, 2021). \emph{Lo que enseña ese nombre no es folclor:} \textbf{enseña que ponerle palabra a lo que se siente es lo que permite pedir ayuda}. \emph{Y enseña algo más, que la investigación documenta y esta guía no va a suavizar:} \textbf{ese sufrimiento venía de las violencias del conflicto armado} ---es decir, mucho de lo que se llama «nervios» tiene una causa afuera, en lo que le pasó a la gente, y no adentro, en su carácter---.}

\begin{drawbox}{milcTurquesa}{La escalera del acercamiento}
\begin{pasoslista}
\item \textbf{Escribe lo que evitas} y ordénalo de lo menos temido a lo más temido.
\item \textbf{Empieza por el primer escalón,} el que te da un 3 o 4 de 10 ---no el que te da 9---.
\item \textbf{Quédate ahí hasta que la sensación baje} un poco, en vez de salir corriendo apenas suba.
\item \textbf{Repite el mismo escalón} varias veces antes de subir. \emph{La repetición es lo que enseña.}
\item \textbf{Si no puedes solo, no es un fracaso:} es la señal de que este acompañamiento le toca a un adulto o a un profesional.
\end{pasoslista}
\end{drawbox}

\begin{softbox}{milcMostaza}{milcGris}{Errores comunes y su corrección}
\textbf{Autodiagnosticarse:} los criterios sirven para decidir si pedir ayuda, no para ponerse una etiqueta. \textbf{«Todo el mundo tiene ansiedad»:} banaliza y deja sola a la persona que sí está incapacitada. \textbf{«Es que eres muy nervioso»:} convierte un problema tratable en un rasgo. \textbf{Evitar y llamarlo cuidarse:} el alivio de hoy es el problema de mañana. \textbf{Lanzarse al escalón más alto:} confirma el miedo en vez de desmentirlo. \textbf{Respirar para escapar:} respirar ayuda a sostener, no a huir. \textbf{Creer que hay que esperar a estar muy mal:} los criterios existen justo para no esperar.
\end{softbox}

\begin{infoband}{milcTurquesa}


\textbf{En tu cuaderno (Actividad 2 · EXPLICA).} Título: «Actividad 2. Los cuatro criterios». \textbf{Formato:} pasa tus cuatro situaciones por interferencia, persistencia, intensidad y malestar, y escribe cuál se parece más a una señal útil y cuál a un problema. \textbf{Extensión:} media página. \textbf{Sabes que terminaste cuando} puedes explicar por qué \textbf{evitar alivia hoy y agrava mañana}.
\end{infoband}

\puente{Ya sabes el mecanismo. Es hora de lo único que lo revierte: acercarse a algo evitado, de a poquitos y con fecha.}

\milcestacion{milcMagenta}{Fase 3 · Praxis: acercarse de a poco}

\begin{softbox}{milcMagenta}{milcGris}{Cómo se deja de evitar sin lanzarse al vacío}
Aquí no se trata de dejar de sentir ansiedad ---eso no se puede, y tampoco convendría---. Se trata de dejar de organizar la vida alrededor de evitarla.
\end{softbox}

\milcpaso{1}{milcMagenta}{\textbf{Las señales del cuerpo (6 min).} Termina la silueta y \textbf{ponle a cada señal una palabra tuya}. \emph{Esto sirve más de lo que parece:} \textbf{poder decir «se me aprieta el pecho y me tiembla la voz» en vez de «me siento mal»} es la diferencia entre que alguien te entienda y que no ---y es exactamente lo que hacía el nombre «de los nervios»---.}
\milcpaso{2}{milcMagenta}{\textbf{Los tres criterios (8 min).} Pasa tus cuatro situaciones por \textbf{interferencia, persistencia e intensidad}, con datos concretos ---horas, días, del 1 al 10---. \emph{No para etiquetarte:} \textbf{para saber si esto lo puedes trabajar tú o si toca contarlo}.}
\milcpaso{3}{milcMagenta}{\textbf{La lista de lo evitado, ordenada (8 min).} Toma tu lista y \textbf{ordénala de menos temido a más temido}, poniéndole a cada cosa un número del 1 al 10. \emph{Casi todo el mundo descubre lo mismo al hacerlo:} \textbf{la lista es más larga de lo que creía}, y varias cosas de la parte baja no son tan graves como se sentían.}
\milcpaso{4}{milcMagenta}{\textbf{El acercamiento pequeño (10 min y la semana).} Escoge \textbf{un solo escalón bajo} ---un 3 o un 4--- y escribe \textbf{qué vas a hacer, qué día y a qué hora}. \emph{Ejemplos reales:} preguntar una cosa en clase; llamar para pedir una cita; saludar a alguien; leer un párrafo en voz alta; ir a un lugar diez minutos. \textbf{Y la regla que hace que funcione:} \emph{quédate hasta que la sensación baje un poco, en vez de irte apenas suba}. \textbf{Después anota qué tan alto llegó y en cuánto tiempo bajó.} \emph{Ese segundo número es el que enseña, porque casi siempre baja antes de lo que uno cree.}}

\begin{drawbox}{milcMagenta}{Antes de cerrar tu mapa}
\begin{checklista}
\item Tus señales del cuerpo tienen \textbf{palabras propias}, no solo «me siento mal».
\item Pasaste las cuatro situaciones por los criterios, con datos concretos.
\item Tu lista de lo evitado está \textbf{ordenada} y numerada del 1 al 10.
\item Escogiste un escalón \textbf{bajo} ---3 o 4---, no el más alto.
\item Tu acercamiento tiene \textbf{día y hora}.
\item Sabes que si no puedes solo, eso \textbf{no es un fracaso}: es cuándo se cuenta.
\end{checklista}
\end{drawbox}

\begin{softbox}{milcMagenta}{milcGris}{Plantilla de trabajo}
\textbf{Nombrar la señal.} \emph{«Se me \_\_\_\_ y \_\_\_\_. Es de \_\_\_\_ de 10 y me dura \_\_\_\_.»} \\[1mm]
\textbf{Los criterios.} \emph{«¿Interfiere? ¿Persiste? ¿Es desproporcionada? ¿Estoy sufriendo?»} \\[1mm]
\textbf{Mi escalón.} \emph{«Voy a \_\_\_\_ el \_\_\_\_ a las \_\_\_\_. Me voy a quedar hasta que baje.»} ---y después: subió a \_\_\_\_, bajó en \_\_\_\_ minutos---.
\end{softbox}

\begin{infoband}{milcMagenta}


\textbf{En tu cuaderno (Actividad 3 · CREA).} Título: «Actividad 3. Mi mapa de la ansiedad». \textbf{Formato:} la silueta con palabras + los criterios aplicados + la lista ordenada y numerada + el acercamiento con día, hora y los dos números del después. \textbf{Extensión:} una página. \textbf{Sabes que terminaste cuando} tienes escrito \textbf{en cuánto tiempo bajó}.
\end{infoband}

\puente{El producto tiene un paso pequeño adentro. \emph{Los grandes no se dan, y los que no se dan no enseñan nada.}}

\milcestacion{milcMostaza}{Producto de la sesión}
\textbf{Mapa de la ansiedad: señales del cuerpo, los criterios aplicados, la lista de lo evitado y un acercamiento con fecha}.
\begin{softbox}{milcMostaza}{milcGris}{Criterios de calidad}
(1) Las señales corporales están nombradas con palabras propias y localizadas. (2) Los criterios de interferencia, persistencia e intensidad se aplican con datos concretos. (3) La lista de lo evitado está ordenada por nivel de temor y numerada. (4) El acercamiento elegido es un escalón bajo, con día y hora, y contempla permanecer hasta que la sensación baje. (5) Se reconoce el límite: cuándo esto deja de manejarse solo y se cuenta a un adulto.
\end{softbox}

\puente{Antes de cerrar, revisa lo importante: si tu acercamiento te da pánico de solo pensarlo, es demasiado grande. Hazlo más pequeño.}

\milcestacion{milcVino}{Evaluación liberadora · cinco dimensiones}

\begin{softbox}{milcVino}{milcGris}{Sentido de la evaluación}
Evaluar no es solo revisar una nota. Es reconocer qué aprendí, qué cambió en mí, qué emociones manejé, cómo me relaciono con la ciudad y la comunidad, y qué le dejaré a quien viene detrás.
\end{softbox}

\textbf{Mi reflexión liberadora en cinco dimensiones.} Completa cada frase iniciada:

\small
\begin{tabularx}{\textwidth}{>{\bfseries\RaggedRight\arraybackslash}p{3.4cm}Y>{\RaggedRight\arraybackslash}p{4.6cm}}
\rowcolor{milcVino}\color{white}Dimensión & \color{white}\bfseries Mi respuesta (frame starter) & \color{white}\bfseries Algo que descubrí\\
1. Personal & Lo que más cambió en mí fue \linead{6.5cm} & \linead{4.2cm}\\
2. Emocional & La emoción más fuerte fue \linead{2.5cm} y la regulé \linead{3cm} & \linead{4.2cm}\\
3. Ciudadana & En democracia esto importa porque \linead{6cm} & \linead{4.2cm}\\
4. Local (Cartago) & En mi barrio sería útil para \linead{6.5cm} & \linead{4.2cm}\\
5. Intergeneracional & A mi abuela/abuelo le diría \linead{6.5cm} & \linead{4.2cm}\\
\end{tabularx}
\normalsize

\begin{infoband}{milcVino}
\textbf{Para la próxima sustentación.} Vuelve a esta tabla en seis meses. Verás que las respuestas cambian — no porque lo aprendido cambie, sino porque tú cambias. La evaluación liberadora es una conversación contigo a través del tiempo.
\end{infoband}


\puente{Cerramos con tres voces: el rostro que reclama, la diferencia entre la cosa y la opinión sobre la cosa, y el cuidado de lo compartido.}

\milcestacion{milcGradeDark}{Triángulo de pensamiento · cierre filosófico}

\begin{softbox}{milcVino}{milcGris}{Dussel · lente del nosotros}
\textit{«El otro se revela realmente como otro\ldots como el pobre, el oprimido; el que, a la vera del camino, fuera del sistema, muestra su rostro sufriente y sin embargo desafiante: «¡Tengo hambre!, ¡tengo derecho a comer!»»}

{\footnotesize\itshape\color{milcNegro!70}--- Enrique Dussel · Filosofía de la liberación (1977)}

La «enfermedad de los nervios» es \textbf{un nombre que se dieron a sí mismas comunidades que estaban fuera del sistema}: sin psicólogo, sin palabra técnica y sin nadie preguntando qué les había pasado. \emph{Y lo que nombraron no era una rareza de su carácter: era lo que les dejó la guerra.} \textbf{Fíjate en lo que eso corrige.} \emph{Cuando el sufrimiento se explica solo por dentro ---«es que es muy nervioso»---, la persona carga con una culpa que no le toca.} \textbf{Preguntar \emph{qué le pasó} en vez de \emph{qué tiene} cambia la conversación entera,} \emph{y esa es la diferencia entre acompañar a alguien y ponerle una etiqueta.}

\textbf{De lo que me angustia, ¿cuánto viene de algo que me pasó y no de cómo soy?}
\end{softbox}
\linead{\linewidth}\\[1mm]
\linead{\linewidth}

\begin{softbox}{milcMostaza}{milcGris}{Estoicismo · Epicteto · Enquiridión, 5 (c. 125 d.C.) · cuidado interior}
\textit{«No son las cosas las que perturban a los hombres, sino las opiniones sobre las cosas.»}

{\footnotesize\itshape\color{milcNegro!70}--- Epicteto · Enquiridión, 5 (c. 125 d.C.)}

\textbf{Cuidado con esta frase, porque mal usada hace daño.} \emph{No significa que si estás angustiado sea culpa tuya por pensar mal ---eso es echarle a la persona la carga de su propio sufrimiento---.} \textbf{Lo que sí dice, y es muy útil, es esto:} \emph{entre lo que pasa y cómo lo vives hay una interpretación, y la interpretación se puede examinar}. \textbf{El caso más claro es el que produce el círculo de la ansiedad:} \emph{el corazón acelerado es una cosa; «me está dando algo grave» es una opinión sobre esa cosa}. \textbf{Y la opinión es la que multiplica el miedo.} \emph{Poder decirse «esto es ansiedad, sube y baja» no quita la sensación ---pero le quita el segundo piso---.}

\textbf{Cuando siento la sensación, ¿qué me digo que ella significa?}
\end{softbox}
\linead{\linewidth}\\[1mm]
\linead{\linewidth}

\begin{softbox}{milcTurquesa}{milcGris}{Floridi · lente de la infoesfera}
\textit{«El florecimiento de las entidades informacionales y de toda la infoesfera debe promoverse preservando, cultivando y enriqueciendo sus propiedades.»}

{\footnotesize\itshape\color{milcNegro!70}--- Luciano Floridi · La ética de la información (2013; trad.)}

Hay una forma de evitación que hoy es la más fácil de todas y no parece evitación: \textbf{la pantalla}. \emph{Cuando algo angustia, abrir el teléfono da alivio inmediato ---y por eso el cerebro aprende a hacerlo---,} \textbf{pero es exactamente el mecanismo que mantiene el problema:} \emph{alivio ahora, y la comprobación de que se podía tolerar nunca ocurre}. \textbf{Y hay una segunda trampa digital que conviene nombrar:} \emph{buscar los síntomas en internet}. \textbf{Casi nunca calma y casi siempre agranda la opinión sobre la cosa} ---justo lo que decía Epicteto---. \emph{Cultivar, aquí, es preferir una conversación con alguien a una búsqueda a las dos de la mañana.}

\textbf{Cuando me angustio, ¿abro el teléfono? ¿Y qué pasó con lo que me angustiaba?}
\end{softbox}
\linead{\linewidth}\\[1mm]
\linead{\linewidth}

\begin{infoband}{milcNegro}
\textbf{Nota para el docente.} Las tres son citas textuales y llevan su referencia debajo. Puedes remitir a la obra en clase.
\end{infoband}

\milcestacion{milcVino}{Mi autoevaluación · marca una casilla por fila}

{\small
\setlength{\extrarowheight}{2.2pt}
\begin{tabularx}{\textwidth}{>{\RaggedRight\arraybackslash\bfseries}p{3.0cm}YYY}
\rowcolor{milcVino}\color{white}\bfseries Criterio & \color{white}\bfseries Lo logré & \color{white}\bfseries En proceso & \color{white}\bfseries Necesito apoyo\\[.6mm]
Comprensión del tema & \milcopcion{Explico las ideas con mis palabras.} & \milcopcion{Reconozco las ideas, falta precisión.} & \milcopcion{Necesito repasar lo básico.}\\[.8mm]
\rowcolor{milcGris}Calidad del mapa & \milcopcion{Distingo la ansiedad que avisa de la que incapacita con criterios, y sé qué he dejado de hacer por evitarla.} & \milcopcion{Reconozco las señales, pero sigo pensando que evitar es la solución.} & \milcopcion{Sigo creyendo que sentir ansiedad ya es estar enfermo.}\\[.8mm]
Aplicación y producto & \milcopcion{Mi producto cumple los criterios.} & \milcopcion{Está armado, con ajustes pendientes.} & \milcopcion{Aún está en construcción.}\\[.8mm]
\rowcolor{milcGris}Reflexión 5 dimensiones & \milcopcion{Respondí cada una con honestidad.} & \milcopcion{Respondí algunas con detalle.} & \milcopcion{Mis respuestas son muy generales.}\\[.8mm]
Triángulo de pensamiento & \milcopcion{Conecté las 3 voces con mi trabajo.} & \milcopcion{Conecté al menos una.} & \milcopcion{No las relaciono con mi proceso.}\\[.8mm]
\end{tabularx}}

\vspace{3mm}
\textbf{Hoy entendí que:}\\[1mm]
\linead{\linewidth}\\[1mm]
\linead{\linewidth}

\textbf{Mi compromiso para la próxima sesión es:}\\[1mm]
\linead{\linewidth}\\[1mm]
\linead{\linewidth}

\begin{softbox}{milcMostaza}{milcGris}{Cita de despedida · Marco Aurelio}
\textit{``El obstáculo en el camino se vuelve el camino. Lo que estorba al camino se vuelve camino.''}\\[1mm]
Cada error de hoy, bien observado, se convierte en pieza del camino que pisarás mañana.
\end{softbox}

\small
\begin{tabularx}{\textwidth}{>{\bfseries}p{3.6cm}Y}
\rowcolor{milcGradeDark!12}Nombre completo: & \linead{10cm}\\
Fecha: & \linead{4cm} \quad Curso/grupo: \linead{4cm}\\
Mi firma: & \linead{9cm}\\
\end{tabularx}
\normalsize

\begin{infoband}{milcGradeDark}
\textbf{Para la próxima sesión.} Dos cosas. \textbf{Primero, da tu escalón} y trae los dos números: \textbf{a cuánto subió y en cuántos minutos bajó}. \emph{Y si la lista de lo que evitas te dejó preocupado ---si son muchas cosas, o si llevan meses, o si te impiden vivir---, cuéntaselo a un adulto esta semana.} \textbf{Los cuatro criterios existen justamente para no esperar a estar muy mal.} \textbf{Segundo, pregunta en tu casa por «los nervios»:} averigua con alguien mayor \textbf{si alguna vez oyó decir que alguien estaba «de los nervios»}, \textbf{qué se decía que había que hacer} ---una agüita, reposo, ir donde alguien, no hablar del tema--- y \textbf{si esa persona había pasado por algo difícil antes}. \textbf{Trae:} tus dos números y lo que te contaron. \emph{Vas a encontrar dos cosas que conviven: un saber muy fino sobre el cuerpo ---la gente sabía exactamente dónde se sentía--- y un silencio muy grande sobre la causa. Casi nunca se preguntaba qué le había pasado a esa persona.}
\end{infoband}

\end{document}
